\subsection{从固定两端的驻波到箱中粒子}
\label{ss:position-well}
长为 $L$、两端固定的弦，其驻波在弦长内恰有整数个半波长。质量为 $m$ 的粒子受限于两面硬墙之间时，两端的边界条件同样限制空间振幅的形状。位置表象使这些条件能够直接写入态方程。
记 $\psi(x,t)=\brkt{x}{\psi(t)}$。对无自旋、无磁场的非相对论粒子，模型取 $\op{H}=\op{P}^2/(2m)+V(\op{X})$。
把式\eqref{eq:schrket}左乘 $\bra{x}$，势能在位置 $x$ 处乘以 $V(x)$，
动量则作用为 $-\i\hbar\p_x$。动量作用两次给出 $-\hbar^2\p_x^2$，于是
\begin{equation}
    \i\hbar\p_t\psi(x,t)
    =-\frac{\hbar^2}{2m}\p_x^2\psi(x,t)+V(x)\psi(x,t).
    \label{eq:position-schrodinger}
\end{equation}
若制备在能量确定为 $E$ 的态，已知它的时间因子是 $\e^{-\i Et/\hbar}$。
代入 $\psi(x,t)=u(x)\e^{-\i Et/\hbar}$，左侧对时间求导得到 $E\psi$，
约去公共的相位，就得到决定空间形状的定态方程
\begin{equation}
    -\frac{\hbar^2}{2m}u''(x)+V(x)u(x)=Eu(x).
    \label{eq:position-stationary}
\end{equation}

取阱内 $0<x<L$ 的势能为零，两端理想硬墙要求 $u(0)=u(L)=0$，阱外振幅为零。
动能非负，而 $E=0$ 的线性解经两端条件只剩零函数。因此令 $k=\sqrt{2mE}/\hbar>0$，
阱内方程就是 $u''=-k^2u$。一般解 $A\sin kx+B\cos kx$ 在左端为零，故 $B=0$。
右端又要求 $\sin(kL)=0$，可归一化的非零解只有
\begin{equation}
    k_n=\frac{n\pi}{L},\qquad
    E_n=\frac{\hbar^2\pi^2n^2}{2mL^2},\qquad
    u_n(x)=\sqrt{\frac{2}{L}}\sin\frac{n\pi x}{L},
    \qquad n=1,2,\ldots
\end{equation}
在这段区间内，正弦平方的平均值为 $1/2$，所以积分为 $L/2$。归一化给出上式的系数，整体相位取零。
和弦的驻波一样，两端条件选出 $L=n\lambda_n/2$。粒子的自由动能为 $\hbar^2k_n^2/(2m)$，
所以第二能级是第一能级的四倍。弦上的纵坐标表示位移，粒子位置探测器读出的是区间里的计数。

\begin{figure}[htbp]
    \centering
    \tikzfigure[width=.85\linewidth]{well-detection}
    \caption{箱中第一激发态 $n=2$。左图给出带符号的空间振幅，右图给出相应的位置概率密度。横坐标为 $x/L$；右图阴影的面积给出探测窗内的概率。}
    \label{fig:well-detection}
\end{figure}
对图\ref{fig:well-detection}中的探测窗 $L/4<x<3L/4$，不同位置的计数相加，所以要对振幅的模平方积分。令 $\xi=x/L$，用 $2\sin^2 y=1-\cos(2y)$，同一扇窗对两个态给出
\begin{equation}
    P_n=\int_{1/4}^{3/4}[1-\cos(2n\pi\xi)]\d\xi,\qquad
    P_1=\frac12+\frac1\pi,\quad P_2=\frac12 .
\end{equation}
若改为测量能量，制备在 $u_2$ 的每个粒子都给出 $E_2$。定态的含义是这些分布不随时间变化。

\subsection{概率怎样流过一个截面}
\label{ss:position-current}
同一批次的粒子穿过一段区域时，区域内的概率会随进出而改变。对实势 $V$，
将式\eqref{eq:position-schrodinger}乘以 $\psi\conj$，再减去它的复共轭，势能项相消，得到
\begin{equation}
    \p_t\abs{\psi}^2
    =\frac{\i\hbar}{2m}\p_x
      \bigp{\psi\conj\p_x\psi-\psi\p_x\psi\conj}
    =-\p_x j,\qquad
    j=\frac{\hbar}{m}\Im(\psi\conj\p_x\psi).
    \label{eq:position-current}
\end{equation}
从 $a$ 积到 $b$，便有 $\d P_{[a,b]}/\nd t=j(a,t)-j(b,t)$。一维的 $j$ 有时间倒数的量纲，
正值表示概率沿 $+x$ 方向流动。箱中的单个定态可取实空间函数，处处有 $j=0$；
行波 $A\e^{\i kx}$ 则给出 $j=(\hbar k/m)\abs{A}^2$，正是速度乘以密度。
后面计算隧穿时，入射与透射的正向流之比就是透射率。关于位置方程、概率守恒及硬墙边界，参见文献\cite{mitwave}。
